\begin{table*}[t]
\centering\footnotesize
\setlength{\tabcolsep}{4pt}
\caption{동일 319장·8코너 평가에서의 보정 방법 비교. 입력·손실·감독이 다른 방법 전체 비교이며 세 seed 통계의 산술평균이다.}\label{tab:comparators}
\begin{tabularx}{\textwidth}{Yrrrrr}
\toprule
방법 & \shortstack{코너 중앙값\\(px)} & \shortstack{코너 P90\\(px)} & \shortstack{\pckten\\(\%)} & \shortstack{위치 중앙값\\(cm)} & \shortstack{회전 중앙값\\(도)} \\
\midrule
Base & 6.721 & 43.890 & 63.425 & 7.897 & 2.539 \\
영상 보정 P & 5.938 & 42.631 & 67.494 & 7.153 & 2.154 \\
직접 회귀 D & 6.504 & 42.992 & 64.372 & 7.597 & 2.274 \\
선 구조 L & 6.146 & 42.984 & 66.867 & 7.548 & 2.233 \\
PoseFix 기반 & 5.561 & 43.907 & 68.707 & 6.951 & 2.028 \\
제안 보정 N3 & 5.778 & 42.134 & 68.587 & 7.068 & 2.070 \\
\bottomrule
\end{tabularx}
\tabnote{Base는 보정 없음, P는 별도 실행의 영상 기반 보정으로 N0와 다른 가중치이다. 같은 평가 집합이 입력·손실·사전학습·선택 예산의 동등성을 의미하지 않는다. PoseFix 기반 방법이 유리한 중앙값을 유지한다.}
\end{table*}
